\todo{Lecture 20}

Dans un diélectrique, les électrons ne peuvent pas bouger librement. Mais un champ $\boldsymbol{E}$ a quand même des effets :
\begin{enumerate}
\item Très faible déplacement des électrons par rapport au noyau des molécules du diélectrique.
\item En présence des moments dipolaires permanents, orientation le long du champ $\boldsymbol{E}$.
\end{enumerate}

La nature d'un diélectrique peut être distingue en deux cas :
\begin{enumerate}
\item Si $\boldsymbol{p}_{\text{permanente}}=\boldsymbol{0}$, une polarisation de la matière est entraînée. Si le diélectrique est homogène, l'effet macroscopique est l'occurrence d'une densité de charge de surface $\sigma _{p}$.
\item Si $\boldsymbol{p}_{\text{permanente}}\neq \boldsymbol{0}$, l'effet net est le même. L'orientation des molécules dépende de la compétition du moment de force $\boldsymbol{M}_{\mathrm{res}}=\boldsymbol{p}\times \boldsymbol{E}$ et l'agitation thermique.
\end{enumerate}
Dans les deux cas, si le diélectrique est homogène (et isotrope) et $\boldsymbol{E}$ n'est pas trop fort, alors
$$
\sigma _{p}\propto E. 
$$
On définit $\mathcal{X}$ tel que $\sigma _{p}=\mathcal{X}\varepsilon_{0}E$, ou $\mathcal{X}$ est la susceptibilité électrique.

On revient au condensateur plan. Par Gauss
\begin{align*}
ES_{c} & =\frac{1}{\varepsilon_{0}}(1-\sigma _{p}S_{c}) \\
 & =\frac{1}{\varepsilon_{0}}(q-S_{c}\mathcal{X}\varepsilon_{0}E)  \\
 \Rightarrow E(\varepsilon_{0}S_{c}+\varepsilon_{0}\mathcal{X}S_{c} ) & =q \\
\Rightarrow E & =\frac{q}{S_{c}\varepsilon_{0}(1+\mathcal{X})}
\end{align*}
donc on observe que $E$ décroît a cause du diélectrique. En particulier
$$
U=Ed=\frac{qd}{S_{c}\varepsilon_{0}(1+\mathcal{X})}
$$
et 
$$
C=\frac{q}{U}=\frac{\varepsilon_{0}(1+\mathcal{X})S_{c}}{d}=\frac{\varepsilon_{0}\varepsilon _{r}S_{c}}{d}
$$
donc la capacite augmente par le facteur $\varepsilon _{r}=(1+\mathcal{X})$ (la permittivité relative).

En general, la capacite d'un conducteur dépende de la géométrie et du diélectrique.  En presence d'un dielectrique $C$ devient $\varepsilon _{r}C$.

\paragraph{Quelques valeures de $\varepsilon _{r}$}
\begin{itemize}
\item Vide : $\varepsilon _{r}=1$
\item Air : $\varepsilon _{r}=1.0006$
\item NaCl : $\varepsilon _{r}=5.6$
\item Eau : $\varepsilon _{r}=8.1$
\item Quelques ceramiques avec $\mathrm{TiO}_{2}$ : $\varepsilon _{r} \underset{\sim}< 5000$
\end{itemize}

\chapter{Circuits electriques}

En electrostatique, pas de courants. Donc $\boldsymbol{E}=\boldsymbol{0}$ a l'interieur des conducteurs. Un appareil a force electromotrice entre ses deux bornes (p.ex une pile) peut mantenir un champ $\boldsymbol{E}$ dans un conducteur. Ceci implique un mouvement des porteurs mobiles de charges (induit un courant)
\subsection{Definition de courant}
\begin{definition}
Le courant $I$ est le flux de charge par unite de temps a travers une surface $S$. Le courant est mesure en amperes $\mathrm{A}=\frac{\mathrm{C}}{\mathrm{s}}$. 
$$
I=n^{+}q^{+}u^{+}S+n^{-}q^{-}u^{-}S
$$
ou $n^{+}$ est la densite des porteurs de charge libres, $q^{+}$ est la charge par particule et $u^{+}$ est la vitesse moyenne des charges.
\end{definition}